\begin{spacing}{1.5}

Small, locally-runnable language models are attractive for education because
they preserve privacy, remove recurring API cost, and work without permanent
dependence on a commercial provider. When used on their own, however, they can
state a wrong or outdated fact with the same fluent confidence as a correct
one, and a learner has no way to tell the difference. This report presents
\textbf{AURA-Bio}, a local-first biology tutoring system built to close that
gap, not by training a bigger model, but by surrounding a small one
(a 4-billion-parameter model served through Ollama) with a pipeline that
routes each question, retrieves supporting evidence from a curated
25\,GiB biology corpus, generates an answer against that evidence, and then
\emph{checks} the answer before it reaches the learner.

The system was built and measured, not just designed. A biology corpus of
25.0\,GiB was assembled from nine official sources (PubMed, PMC, Wikipedia,
UniProt, Gene Ontology, MeSH, Reactome, NLM LitArch, and open textbooks),
verified file-by-file, deduplicated, and chunked into over ten million
passages, of which 10,394,250 are embedded for dense retrieval and combined
with a BM25 sparse index through reciprocal rank fusion and cross-encoder
reranking. A rule-and-classifier router assigns each question a task type,
biology subdomain, and learner level. Generated answers pass through a
verification layer that recomputes rather than merely pattern-matches:
SymPy re-derives genetics and population-genetics arithmetic, Biopython
checks codon and sequence claims against the standard genetic code,
\texttt{periodictable} balances stated chemical equations by element, and a
natural-language-inference model checks whether a claim actually agrees with
its cited evidence rather than merely sharing its vocabulary. A curated
lexicon of 135,941 entity names drawn from the corpus's own authorities
(Gene Ontology, MeSH, Reactome, and the encyclopaedia index) lets the system
recognise when a question is built on something that does not exist. An
external critic model reviews answers during development, and every stage of
the pipeline is traced for later inspection.

A 441-question benchmark, built with machine-checkable answer keys rather
than just another second language model as judge, was used throughout to
measure, not assume, whether each addition helped. That discipline
produced the most important finding of the internship: the underlying
4-billion-parameter model already answers the great majority of these
questions correctly on its own, so raw accuracy leaves little room for any
addition to move the needle, and the pipeline's real contribution lies
elsewhere, in making every claim traceable, checkable, and correctly
refused when the premise is false.

The measured results appear in Chapter~4 (overall accuracy \textbf{0.889},
statistically indistinguishable from the unaided baseline at this sample
size,
retrieval's headroom bounded to \textbf{six percentage points} by how much
of the benchmark the raw model already knows, \textbf{60.5\%} of the
pipeline's claims carrying a verifiable citation against \textbf{zero} for
the unaided model, and a hallucination rate of \textbf{12.5\%} on
deliberately adversarial questions for both systems). Fine-tuning the base
model on a curated instruction set and prioritising explicit refusal
examples over raw additional text, for reasons argued from the corpus's own
token economics is scoped in detail and is the immediate next phase of the
project. The
report closes with a plan for generalising the same architecture beyond
biology to other STEM domains through a modular plugin interface, so that
the retrieval, verification, and critique machinery built here can be reused
rather than rebuilt for mathematics, physics, or coding.

\end{spacing}
